\begin{table}[h]
\centering
\caption{Replication of Table 1, Panels A \& B: Effect of AA on SAT Scores}
\label{tab:table1_ab}
\small
\begin{tabular}{lcc}
\toprule
 & Math & Verbal \\
 & (1)  & (2)    \\
\midrule
\multicolumn{3}{l}{\textit{Panel A: URMs (Black and Hispanic)}} \\[4pt]
DD coefficient & 8.009*** & -0.634 \\
               & (1.565)   & (1.808)   \\
\\[-2pt]
Observations   & 1,904 & 1,901 \\
$R^2$          & 0.844 & 0.796 \\
State, Year, Race FE & X & X \\
\midrule
\multicolumn{3}{l}{\textit{Panel B: Whites}} \\[4pt]
DD coefficient & 4.048*** & 0.034 \\
               & (1.024)   & (0.925)   \\
\\[-2pt]
Observations   & 663 & 663 \\
$R^2$          & 0.969 & 0.971 \\
State, Year, Race FE & X & X \\
\bottomrule
\end{tabular}
\begin{minipage}{\linewidth}
\smallskip
\footnotesize
\textit{Notes}: Difference-in-differences estimates of the effect of reinstating
affirmative action (AA) on mean SAT scores at the state--race--year level.
The DD coefficient is the interaction of an indicator for a treated state
(Texas, Louisiana, Mississippi) and an indicator for post-2003 (post-\textit{Grutter v.\
Bollinger}). Regressions include controls for AA-ban policy changes in
control states. Cells are weighted by the number of test-takers.
Standard errors (in parentheses) are clustered at the state level.
$^{*}p<0.10$, $^{**}p<0.05$, $^{***}p<0.01$.
\end{minipage}
\end{table}
